% Nouveau EGA V (cote 25): archive cote 25, site page 33 (PDF page 34).
% typescript. Draft compared: none.
% AI-transcribed by Claude Opus 5.5 (claude-opus-5-5, adaptive thinking, effort high), Message Batches API, batch ega25b (msgbatch_01A7n5YuXubUPNK1zsepi8RG), 2026-11-13; original AI transcription, subsequently checked against both supplied scan parts; see PAGES_CHECKED.tsv and CHECK_DECISIONS.tsv.
\scanpage{25/33}
\mspage{\hand{29}}

Noter qu'avec les notations précédentes, si $\hand{\xi\in}$ $\overset{\hand{\vee}}{P}$ - Z , alors pour tout
y $\hand{\in}$ Y$_{\hand{\xi}}$ , on a coprof$_y$ F$_{k(\hand{\xi})}$ = coprof$_y$ G$_{\hand{\xi}}$ , et par suite: si coprof F$\hand{\leq}$n,
alors pour $\hand{\xi\in}\overset{\hand{\vee}}{P}$ - Z, on a coprof G$_{\hand{\xi}}$ $\hand{\leq}$ n; en particulier si F est CM, alors
pour $\hand{\xi\in}$ $\overset{\hand{\vee}}{P}$ - Z, G$_{\hand{\xi}}$ est CM. \var{\inserted{Enfin, si F est $(S_k)$, on sait que $G_\xi$ est $(S_{k-1})$ pour $\xi\in\overset{\hand{\vee}}{P}-Z$ (réf $E_{\uline{\mathrm{IV}}}$).}}{Enfin, si F est $(S_k)$, on sait que $G_\xi$ est $(S_{k-1})$ pour $\xi\in\overset{\hand{\vee}}{P}-Z$ (réf $E_{\uline{\mathrm{IV}}}$).}{insertion}{ega25_25_033-v0001}
\hand{8.3.} On notera que \var{\deleted{on n'a pas de résultat analogue pour la propriété}}{}{deletion}{ega25_25_033-v0002} \var{\inserted{\illegible{4} en général}}{\illegible{4} en général}{insertion}{ega25_25_033-v0003}
\var{\deleted{(S$_k$)}}{}{deletion}{ega25_25_033-v0004}, même si F=$\underline{O}_X$ , k=1, X étant géom intègre de dim.2 (ou \var{\inserted[typed]{k=2,}}{k=2,}{insertion}{ega25_25_033-v0005} X étant
géométriquement intègre et géométriquement normal de dim.3). Il suffit
de partir d'une surface projective \var{\xout{xxxxxx}}{}{typewriter strike-out}{ega25_25_033-v0006} intègre \var{\inserted[typed]{X $\hand{\subset}$ P$^r$}}{X $\hand{\subset}$ P$^r$}{insertion}{ega25_25_033-v0007} sur k algébriquement
clos, ayant un point \var{\inserted[typed]{x}}{x}{insertion}{ega25_25_033-v0008} où X n'est pas Cohen-Macaulay; alors pour tout
hyperplan passant par x, la section hyperplane Y$_{\hand{\xi}}$ correspondante admet
x pour cycle associé immergé (resp. on part d'une variété normale\hand{,} \hand{(}donc
\var{\inserted[typed]{(S$_2$)\hand{)},}}{(S$_2$)\hand{)},}{insertion}{ega25_25_033-v0009} intègre X $\hand{\subset}$ P$^r$ de dimension 3, ayant un point x \hand{$\in$} X en lequel X n'est pas
Cohen-Maaaulay, alors les Y passant par x ne sont pas \var{\deleted{\unsure{(S}}}{}{deletion}{ega25_25_033-v0010} CM i.e. ne sont
pas (S$_2$) en x)\hand{.} \var{\over{d}{D}}{D}{writing over}{ega25_25_033-v0011}ans ces exemples, l'ensemble des \hand{$\xi$} ``exceptionnels''
pour la propriété (S$_k$) contieht l'hyperplan d e $\overset{\hand{\vee}}{P}$ défini par x\hand{$\in$}P, et est
de codimension 1 (et non de codimension$\hand{\geq}$2). \var{\inserted{Cf \var{\deleted{\unsure{plus haut}}}{}{deletion}{ega25_25_033-v0013} 8.5. plus bas pour \var{\deleted{un}}{}{deletion}{ega25_25_033-v0014} résultat général précis dans cette voie.}}{Cf \var{\deleted{\unsure{plus haut}}}{}{deletion}{ega25_25_033-v0013} 8.5. plus bas pour \var{\deleted{un}}{}{deletion}{ega25_25_033-v0014} résultat général précis dans cette voie.}{insertion}{ega25_25_033-v0012}

\var{\deleted{Cependant, si F satisfait à la propriété (S$_k$), on peut \unsure{montrer|monteer} à
l'aide de ce qui précède que l'ensemble des $\hand{\xi\in}$ $\overset{\hand{\vee}}{P}$ tels que $\phi_{\hand{\xi}}$ soit
F$_{k(\hand{\xi})}$-régulier\hand{,} \uline{et} G$_{\hand{\xi}}$ satisfai\var{\xout{\illegible{1}}}{}{typewriter strike-out}{ega25_25_033-v0016}t (S$_{k-1}$), a un complémentaire qui est de
codimmnsion $\hand{\geq}$ 2. Cela résulte du critère 5.7.4. via la coppofondeur,
du fait que l'ensemble Z$_n$ des points de X où F est d e coprofondeur $>$n est
fermé, et que la codimension de Z$_{n\hand{\xi}}$ dans X$_{\hand{\xi}}$= Y$_{\hand{\xi}}$ , même si elle peut
faire (pour d es \hand{$\xi$} exceptionnels) des sauts \var{\deleted{en hauteur}}{}{deletion}{ega25_25_033-v0017} \var{\inserted{vers le bas}}{vers le bas}{insertion}{ega25_25_033-v0018}, ledit saut ne peut
jamais être que d'une unité au plus, pour des raisons de dimension éviden-%
tes.}}{}{deletion}{ega25_25_033-v0015}

\var{\inserted{Si F est $(S_k)$, que l'ens. des $\xi\in\overset{\hand{\vee}}{P}$ tels que $\phi_\xi$ soit $F_{k(\xi)}$-régulier et $G_\xi$ soit $(S_k)$ ait un complémentaire de codim $\hand{\geq} 2$,}}{Si F est $(S_k)$, que l'ens. des $\xi\in\overset{\hand{\vee}}{P}$ tels que $\phi_\xi$ soit $F_{k(\xi)}$-régulier et $G_\xi$ soit $(S_k)$ ait un complémentaire de codim $\hand{\geq} 2$,}{insertion}{ega25_25_033-v0019}

\var{\margin{33}}{\placedmargin{33}}{marginal note placed}{ega25_25_033-v0020}
